\section{Conclusion}\label{conclusion:conclusion}

We tested four client-centric consistency properties in a three-node Cassandra cluster by varying consistency levels, inter-node delay, network availability, and write timestamps. Under weak read and write settings, delayed propagation exposed read-your-writes and monotonic-reads violations when the client accessed different coordinators. QUORUM/QUORUM and ALL/ONE produced no observed violations in the corresponding successful read experiments. ONE/QUORUM gave mixed results: RYW violations occurred in some runs, while no MR violations were observed under the tested paths and schedules. The latter result does not establish a general MR guarantee.

The monotonic-writes experiment showed that natural timestamp ordering preserved the intended final overwrite in all successful trials. When we deliberately assigned the earlier write a greater timestamp, that write won every successful verification read, regardless of the write consistency level. This measures the final reconciled value, rather than the order in which each replica applied the writes. In the writes-follow-reads experiment, no loss of the dependent write was observed in 1,896 natural-mode trials that read the target version. With inverted client timestamps, the dependent write lost in all 2,200 triggered trials, including those using stronger consistency levels.

Failures and partitions also changed which operations could complete. A majority partition could still serve quorum requests, whereas the isolated node could complete only operations that required a single replica. Taken together, these results show that consistency levels affect visibility and availability, while timestamp ordering determines the winner of competing writes in our MW and WFR tests. The findings describe the observed workloads and success conditions; they do not constitute a formal verification of the four consistency properties.
